\section{Equilibrium and Benchmarks}\label{sec:equilibrium}

This section reduces the participation subgame to a smooth system and defines the two nested environments used to identify the role of private repricing.

\subsection{Participation fixed point}

Let
\[
\delta=\alpha\beta,\quad
A_i=v+\alpha(\rho+x_i),\quad
A_T=v+\alpha\rho_T,\quad
d=\kappa_L-\kappa_T-p_T .
\]
In the central regime, local project mass is $n_i^F=1-t_i$ and private project mass is
\[
n_T^F=t_1+t_2-\frac{2(\kappa_T+p_T)}{q_T}.
\]
Using \eqref{eq:partnerlocal}--\eqref{eq:partnerprivate}, the project-quality indices satisfy
\begin{align}
q_i &= A_i+\delta(1-t_i), \label{eq:qi}\\
q_T &= A_T+\delta\left(t_1+t_2-\frac{2(\kappa_T+p_T)}{q_T}\right). \label{eq:qT}
\end{align}
Together with the indifference condition $t_i(q_i-q_T)=d$, the participation equilibrium solves
\begin{align}
F_i &\equiv t_i\left[A_i+\delta(1-t_i)-q_T\right]-d=0,\qquad i=1,2, \label{eq:Fi}\\
F_T &\equiv q_T-A_T-\delta\left[t_1+t_2-\frac{2(\kappa_T+p_T)}{q_T}\right]=0. \label{eq:FT}
\end{align}
This three-equation system is the smooth object used for implicit differentiation and numerical verification within the central interior regime.

\subsection{Private pricing}

Let $D_T\equiv n_T^F$ denote private-route project demand after participation has adjusted. An interior private price satisfies
\begin{equation}
D_T+p_TD_{T,p}=0,
\label{eq:privatefoc}
\end{equation}
with second-order condition
\begin{equation}
2D_{T,p}+p_TD_{T,pp}<0.
\label{eq:privatesoc}
\end{equation}
Implicit differentiation of \eqref{eq:privatefoc} yields the policy response
\begin{equation}
p^*_{T,x_i}
=
-\frac{D_{T,x_i}+p_TD_{T,px_i}}
{2D_{T,p}+p_TD_{T,pp}}.
\label{eq:px}
\end{equation}
At the reported full-game stationary candidate this derivative is negative. Economically, stronger regional facilitation draws projects away from the shared private route; the private intermediary responds by lowering its optimal access price.

The negative price response is not itself a new mechanism. Indeed, when $\beta=0$ the private demand problem becomes linear and the sign of the price response can be established directly. The role of two-sidedness in our main result is different: it creates the positive public--public cross effect in the fixed-price benchmark.

\subsection{Nested benchmarks}

We compare two environments with identical primitives.

\paragraph{Fixed-price benchmark (B3).}
Both regional public hubs and the shared private hub operate, and participation remains two-sided, but the private access price is held fixed while a public deviation is evaluated. Along the G stationary branch used by the local theorem, define the matched scalar fee
\begin{equation}
\bar p_T(\beta)
= p_T^*\!\left(x_1^G(\beta),x_2^G(\beta);\beta\right).
\label{eq:matchedpricepath}
\end{equation}
For each fixed $\beta$, every B3 public deviation is evaluated holding this scalar $\bar p_T(\beta)$ fixed. The public choices and resulting allocation are otherwise allowed to readjust within B3. Hence B3 and G use the same matched scalar fee for the comparison, but their stationary public investments generally differ: $x_i^{B3}(\beta)$ need not equal $x_i^G(\beta)$.

This distinction also separates two derivative objects. The chain-rule decomposition of the full-game objective can switch the price-policy derivative off locally at a given G state. The headline B3--G comparison instead compares the local strategic slopes at the respective B3 and G stationary states associated with the same matched fee path. Those are not the same-state allocation.

B3 is therefore an identification benchmark for the private repricing margin. It suppresses the policy response $p_T^*(x_1,x_2)$ while avoiding a mechanical difference in the matched fee level. It is not a planner problem, a literal price-regulation counterfactual, or a claim that public investments are complements for every exogenous value of $\bar p_T$.

\paragraph{Full game (G).}
The private intermediary observes $(x_1,x_2)$ and optimally chooses $p_T$ according to \eqref{eq:privatefoc}. Governments anticipate the reduced price policy $p_T^*(x_1,x_2)$ when choosing facilitation investments.

Let $\widetilde W_i$ denote government $i$'s reduced objective after downstream private pricing and participation have been solved. At an interior local optimum,
\begin{equation}
\BR_i'(x_j)
=
-\frac{\widetilde W_{i,x_ix_j}}
{\widetilde W_{i,x_ix_i}}.
\label{eq:brslope}
\end{equation}
Because the own second derivative is negative at the regular interior configurations considered below, the sign of the best-response slope is the sign of the public cross derivative.
